\section{پاسخ سوال ۱}
\begin{stepbox}
\field{تحلیل پدیده کاهش سطح دریچه (Aperture) در دوربین Pinhole}{
در دوربین‌های Pinhole (و به طور کلی تمامی دوربین‌ها)، تغییر اندازه دریچه یا Aperture اثرات فیزیکی و اپتیکی مستقیمی بر روی تصویر ثبت‌شده دارد. با \textbf{کاهش سطح دریچه}، دو پدیده متضاد (یکی مثبت و دیگری منفی) رخ می‌دهد:

\begin{itemize}
    \item \textbf{پدیده مثبت (افزایش عمق میدان و وضوح):} با کوچک‌تر شدن دریچه، پرتوهای نوری که از یک نقطه از جسم ساطع می‌شوند، با زاویه بسیار محدودتری وارد دوربین می‌شوند. این امر باعث می‌شود که دایره اغتشاش (Circle of Confusion) روی سنسور به حداقل برسد. در نتیجه، اشیاء در فواصل مختلف به طور همزمان در فوکوس قرار می‌گیرند (عمق میدان یا Depth of Field افزایش می‌یابد) و تصویر واضح‌تر و شارپ‌تر می‌شود.
    
    \item \textbf{پدیده منفی (کاهش نور ورودی و پدیده پراش):} کوچک شدن دریچه به معنای کاهش سطح مقطع عبور نور است. بنابراین، فوتون‌های کمتری به سنسور می‌رسند که منجر به \textbf{تاریک شدن تصویر} (Underexposure) می‌گردد. برای جبران این کمبود نور، باید زمان نوردهی (Shutter Speed) یا حساسیت (ISO) را افزایش داد که به ترتیب موجب تاری حرکتی (Motion Blur) و افزایش نویز (Noise) می‌شود. علاوه بر این، در دریچه‌های بسیار کوچک، پدیده فیزیکی \textbf{پراش (Diffraction)} نور در لبه‌های روزنه رخ می‌دهد که باعث پخش شدن نور و کاهش وضوح کلی تصویر (Softness) می‌گردد.
\end{itemize}
}
\end{stepbox}
